\uuid{eTmW}
\exo7id{7475}
\titre{exo7 7475}
\auteur{exo7}
\organisation{exo7}
\datecreate{2021-08-10}
\isIndication{false}
\isCorrection{false}
\chapitre{Géométrie affine euclidienne}
\sousChapitre{Géométrie affine euclidienne du plan}
\cpge{mpsi-mp2i.sous-espaces-affines, ptsi.espaces-vectoriels.b, mpsi-mp2i.applications-lineaires.b, pcsi.applications-lineaires.b, ptsi.applications-lineaires.b, mp-mpi.endomorphismes-euclidien.c, pc.endomorphismes-euclidiens.a, psi.endomorphismes-euclidien.a, pt.isometries-euclidien.a, mp-mpi.endomorphismes-euclidien.d, pc.endomorphismes-euclidiens.c, psi.endomorphismes-euclidien.d, pt.isometries-euclidien.c}
\module{Géométrie}
\niveau{L2}
\difficulte{}

\contenu{
\texte{
Soit $E$ un plan affine euclidien. Soient $\mathcal{R}$ un repère 
cartésien orthonormé de $E$,
et $f : E\to E$ l'application affine
définie dans $\mathcal{R}$ par l'égalité
$$f \bigl((x,y) \bigr)=\biggl({3x+4y+8\over 5},{4x-3y-1\over 5}
\biggr).$$
}
\begin{enumerate}
    \item \question{L'application $f$ possède-t-elle des points fixes ?}
    \item \question{Démontrer que lorsque $M$ décrit $E$, le milieu de 
$(M,f(M))$ décrit une droite dont on précisera
l'équation.}
    \item \question{Démontrer que $f$ est une isométrie dont on précisera la 
nature, l'axe et la composante translation.}
\end{enumerate}
}
